\documentclass[../dayan_kernel_manual.tex]{subfiles}

\begin{document}

\section{Host RPC 与异常 ABI}\label{sec:kernel-060}

本章描述 Runtime 与 Host 的内部协议。算子通过 \code{printf}、allocator 和 DAPI
间接使用这些机制，但 \code{rpc_call}、mailbox MMIO 和 trap 结构均不是算子可链接
API。

\subsection{64-bit command word}

每条 RPC 命令是一个 64-bit word：

\begin{center}
\code{command[63:48] = opcode, command[47:0] = arg}
\end{center}

opcode 为 16 bit；arg 通常是 Device 地址或由服务定义的立即数。Runtime 在打包指针
时检查地址高 16 位为 0，Host 不得擅自符号扩展或重解释低 48 位。

\subsection{mailbox 与 doorbell}

下表地址均是相对 \code{CONFIG_CONTROL_BASE} 的 offset，默认 control base 为 0。

\begin{dayantable}{RPC 硬件寄存器}{tab:kernel-060}{X[30,l] X[25,l] X[45,l]}
  寄存器 & offset & 用途 \\
  Host--Device mailbox & \code{0x1000 + hart*8} & Host 写入命令；也是 BL0 使用的每核启动 mailbox。 \\
  \code{IRQ_STATUS} & \code{0x2000} & Device 轮询自己的 outbound pending bit；Host 清除后可发送下一笔。 \\
  \code{IRQ_STATUS_SET} & \code{0x2008} & Device 写 \code{1 << hart} 通知 Host。 \\
  \code{IRQ_STATUS_CLR} & \code{0x200c} & 初始化时清除本 hart 遗留 outbound 状态。 \\
  \code{IRQ_ENABLE_SET} & \code{0x2010} & 初始化时使能本 hart 对应的 Host IRQ 输出。 \\
  Device--Host mailbox & \code{0x2020 + hart*8} & Device 写入 outbound command。 \\
  CLINT \code{msip} & \code{0x20000 + hart*4} & Host 写 1 通知 Device；Device 完成接收处理后写 0。 \\
\end{dayantable}

发送方必须先写 mailbox，再置 doorbell；接收方读取并消费 command 后清除 doorbell。
发送方在前一 doorbell 清除前不得复用 mailbox。MMIO 操作两侧使用
\code{fence iorw,iorw} 保证顺序。

\subsection{每核双 RPC proc}

每个 hart 在启动时创建 receive 和 send 两个内部 proc，并各自使用链接器预留的
4 KiB 栈。

\subsubsection{receive proc}

\begin{enumerate}
  \item 用 Runtime poll wait 检查本 hart \code{msip}；
  \item 读取 Host--Device mailbox 并拆分 opcode/arg；
  \item 调用弱的 \code{rpc_receive(opcode,arg,response_command)} 分派；
  \item handler 返回正数时，把完整 response command 压入本核 send queue；返回 0
        表示无需响应，负数表示请求已消费但失败；
  \item 写 0 清除 \code{msip}，允许 Host 复用 inbound mailbox；
  \item 需要响应时等待该发送事务完成，再处理下一条 inbound command。
\end{enumerate}

handler 在清除当前 \code{msip} 之前运行，不应调用需要同一 Host 线程嵌套处理的阻塞
反向 RPC，否则 Host 与 Device 可能互相等待。

\subsubsection{send proc}

send proc 独占本 hart 的 intrusive transaction queue。队列为空时挂入事件链表；
提交者入队后唤醒它。发送流程为：

\begin{enumerate}
  \item poll \code{IRQ_STATUS}，确认本 hart bit 已清除；
  \item 写 Device--Host mailbox，再写 \code{IRQ_STATUS_SET}；
  \item 再次以 poll wait 等待 Host 清除该 bit；
  \item 调用 transaction cleanup，标记完成并唤醒全部本地等待者。
\end{enumerate}

普通算子调用 \code{rpc_call} 时，request、response 与 packet 在调用 proc 的栈/内存中
保持有效，直至 send proc 收到 Host 完成确认并唤醒调用者。

\subsection{Device 请求 packet}

除异常报告外，Device 主动请求的 arg 指向 40-byte \code{struct rpc_packet}：

\begin{dayantable}{\code{struct rpc_packet} ABI}{tab:kernel-061}{X[17,l] X[28,l] X[15,l] X[40,l]}
  offset & 字段 & 类型 & 含义 \\
  0 & \code{request_addr} & \code{uint64_t} & 请求数据地址；长度为 0 时可以为 0。 \\
  8 & \code{request_size} & \code{uint64_t} & 请求字节数。 \\
  16 & \code{response_addr} & \code{uint64_t} & Host 写响应的 Device 地址。 \\
  24 & \code{response_size} & \code{uint64_t} & 响应区容量/约定长度。 \\
  32 & \code{status} & \code{int64_t} & Host 回填；非负表示成功，负值表示失败。 \\
\end{dayantable}

\begin{dayantable}{当前 Device--Host service opcode}{tab:kernel-062}{X[18,l] X[32,l] X[50,l]}
  opcode & 名称 & arg/packet 语义 \\
  1 & \code{RPC_WRITE} & packet request 指向 24-byte print request，包含 hart/proc ID 和待输出数据地址/长度，无响应区。 \\
  2 & \code{RPC_CLEAR} & packet request 同样使用 print request，其中数据地址/长度为 0，清空对应 hart/proc 的 Host 输出缓冲。 \\
  3 & \code{RPC_HEAP_ACQUIRE} & request 为一个 64-bit size；response 为 64-bit addr 与 64-bit size。 \\
  4 & \code{RPC_HEAP_RELEASE} & request 为 64-bit addr 与 64-bit size，无响应区。 \\
  5 & \code{RPC_TRAP_REPORT} & 特例：arg 直接指向 320-byte trap record，不使用 rpc packet。 \\
\end{dayantable}

\code{rpc_print_request} 的 offset 0/4 是 32-bit \code{hart_id}/\code{proc_id}，
offset 8/16 是 64-bit \code{data_addr}/\code{data_size}，总大小 24 bytes。Host 可以
据此为不同 hart 和 proc 分离输出流；WRITE 的实际字符串位于 \code{data_addr}，不是
紧随 rpc packet 或 print request 存放。

Device 在发起请求前 clean request、response 和 packet 所在 cache line。WRITE 还会
单独 clean 输出数据区。Host 完成后，
Device invalidate packet；status 非负时再 invalidate response。Host 必须保证在清除
outbound interrupt bit 前已经写完 status/response，并且不会在清除后继续修改它们。

\subsection{功能中断与 trap}

Dayan 是非抢占式 Runtime。BL1 设置 \code{mie=0} 并清除 \code{mstatus.MIE}，因此
Host 写 \code{msip} 只改变可轮询状态，不会让 receive proc 抢占普通算子。

同步异常仍通过 \code{mtvec} 进入 trap。BL1 把每 hart trap 栈顶写入
\code{mscratch}；汇编入口用 \code{csrrw sp,mscratch,sp} 切换到独立 trap 栈，保存
x0--x31（x2 记录异常前 sp）、\code{mstatus}、\code{mepc}、\code{mcause}、
\code{mtval} 和 \code{fcsr}。

C handler 生成 Host 可读的稳定记录：

\begin{dayantable}{\code{struct trap_record} ABI}{tab:kernel-063}{X[17,l] X[26,l] X[17,l] X[40,l]}
  offset & 字段 & 类型 & 含义 \\
  0 & \code{version} & \code{uint32_t} & 当前为 1。 \\
  4 & \code{size} & \code{uint32_t} & 当前为 320。 \\
  8 & \code{hart_id} & \code{uint32_t} & 发生异常的 hart。 \\
  12 & \code{thread_id} & \code{uint32_t} & 当前 thread ID。 \\
  16 & \code{thread_type} & \code{uint32_t} & BASE=0，PROC=1。 \\
  20 & \code{thread_state} & \code{uint32_t} & 异常时状态，通常为 RUNNING=0。 \\
  24--56 & machine/F CSR & 5 $\times$ \code{uint64_t} & 依次为 mstatus、mepc、mcause、mtval、fcsr。 \\
  64--319 & \code{x[32]} & 32 $\times$ \code{uint64_t} & x0--x31 完整整数寄存器现场。 \\
\end{dayantable}

记录被复制到每核 RPC 状态中的稳定槽并 clean cache，然后作为 opcode 5 事务异步排入
send queue。普通故障 proc 随即进入 EXIT 状态并切回 scheduler；Host 确认前槽不会
复用。若报告无法入队、RPC send proc 自身损坏、发生嵌套异常或异常来自不能安全退出的
init/scheduler 上下文，Runtime 停机而不是静默丢失诊断。

\subsection{Host 实现检查点}

\begin{itemize}
  \item 每 hart inbound/outbound mailbox 各只允许一笔未确认事务；
  \item 严格遵守 mailbox-before-doorbell 和 data/status-before-clear 顺序；
  \item 正确访问 Device 地址，并与 Runtime 的 clean/invalidate 约定匹配；
  \item 对 packet 和 trap record 使用本章给出的固定宽度与 offset，不使用 Host C
        编译器的隐式结构体布局猜测；
  \item heap acquire 返回的 region 必须可由对应 hart 普通访问，且生命周期持续到
        release 成功；返回大小不得小于请求值；
  \item heap region 的 cache ownership 遵循 Device 接入前 invalidate、归还前 clean、
        Host 确认后 invalidate 的顺序；
  \item shared coherent allocation 必须位于 \code{0x00400000--0x007fffff}，且在
        所有使用者结束前不得回收。
\end{itemize}

\end{document}
